%%%%%%%%%%%%%%%%%%%%%%%%%%%%%%%%%%%%%%%%%%%%%%%%%%%%%%%%%%%%%%%%%%%%%%%%%%%%%%%%
%%  EINLAGE 5: IHR VORTRAG ÜBER IHR PROJEKT
%%  Folien-Upload: 27.01.2027 | Vortrag: 28.01.2027 | vorgeschlagene Abgabe: 07.02.2027 (freiwillig)
%%  Struktur: Beschreibung (1/4–1/2 S.) + Reflexion (1–2 S.)
%%  Einlage = die Folien des Teils, den DU gehalten hast -> Anhang.
%%%%%%%%%%%%%%%%%%%%%%%%%%%%%%%%%%%%%%%%%%%%%%%%%%%%%%%%%%%%%%%%%%%%%%%%%%%%%%%%
\section{Einlage: Ihr Vortrag über ihr Projekt.}
\label{sec:vortrag}

\begin{hinweise}
Hier nehmen Sie die von Ihnen erstellten und präsentierten Folien des Projektvortrages und diskutieren sie
(nur den Teil, den Sie gehalten haben). Für die Moodle-Abgabe am 07.02.: \enquote{Die Folien ihres Vortrages
mit Kommentar und Reflexion dazu}.
\end{hinweise}

\subsection{Beschreibung}
\begin{hinweise}
Die Beschreibung gibt alle notwendigen Informationen an, die die Einlage beschreiben, also zum Beispiel das
Datum, die Art der Folien, die verwendete Software und eine kurze Inhaltsübersicht, aber keine Nacherzählung
(min. $1/4$ Seite, max. $1/2$ Seite!).
\end{hinweise}
\begin{umfang}{Einlage 5 (Vortrag): Beschreibung}{0.25}{0.5}
\platzhalter{Datum (\Vortragsdatum), Anlass, Art der Folien und Software (z.\,B. LaTeX-Beamer mit der
Folien-Vorlage), Anzahl Folien, eigener Anteil vs.\ \Partner, kurze Inhaltsübersicht; Folien in
Anhang~\ref{anh:vortrag}}
\end{umfang}

\subsection{Reflexion}
\begin{hinweise}
In der Reflexion sollen Sie sich Gedanken machen über die Wirkung dieser Einlage und der Arbeit mit ihr auf
Ihren Lernfortschritt.
\begin{itemize}\itemsep0pt
\item In welchem Zusammenhang steht die Einlage mit Ihrem Projektthema?
\item In welchem Zusammenhang steht die Einlage mit Ihrem Schwerpunktthema?
\item Was haben Sie bei der Erstellung dieser Einlage gelernt?
\item Wie bewerten sie die Zusammenarbeit in Ihrer Gruppe?
\item Wie hat Ihnen die Einlage beim Verständnis Ihres Projektthemas geholfen?
\item Wie hat ihre \textbf{wissenschaftliche Arbeit} zur Argumentation beigetragen?
\item Wie hat Ihnen die Einlage beim Verständnis Ihres Schwerpunktthemas geholfen?
\item Wo gab es Schwierigkeiten bei der Beschäftigung mit der Einlage?
\item Was hätten Sie anders gemacht, wenn Sie die Einlage nochmals erstellen müssten?
\item Was sagen Ihre Kommilitoninnen oder Kommilitonen dazu, als Kommentar oder wenn Sie mit Ihnen über die Einlage diskutieren?
\item Enthält Ihre Präsentation logische oder physikalische Fehler?
\item Gab es bei der Beschäftigung mit der Einlage eine neue Erkenntnis?
\end{itemize}
\textbf{Die Reflexion für eine Einlage sollte mindestens eine Seite und maximal 2 Seiten lang sein. Die
Längenangaben sind ohne Titel, Bilder, Tabellen und Tabellen- oder Bildlegenden zu verstehen.}
Tipp: Fragen/Feedback aus der Diskussion nach dem Vortrag notieren und zitieren.
\end{hinweise}
\begin{umfang}{Einlage 5 (Vortrag): Reflexion}{1}{2}
\platzhalter{Reflexion in Fließtext -- Leitfragen erkennbar beantworten; einzelne Folien konkret
ansprechen (\enquote{Folie 4: \dots})}
\end{umfang}

%\subsection{Eigene Bewertungskriterien}
%...
